\documentclass[11pt,a4paper]{article}
\usepackage[spanish]{babel}
\usepackage[utf8]{inputenc}
\usepackage[T1]{fontenc}
\usepackage{amsmath,amssymb,bm}
\usepackage{geometry}
\usepackage{xcolor}
\usepackage{booktabs}
\usepackage{hyperref}
\geometry{margin=2.5cm}
\hypersetup{colorlinks=true,linkcolor=blue,urlcolor=blue}

\newcommand{\BHF}{B_{\mathrm{HF}}}
\newcommand{\DEQ}{\Delta E_Q}
\newcommand{\dd}{\mathrm{d}}

\title{Distribuciones bidimensionales $P(\BHF,\DEQ)$ en espectros Mössbauer}
\author{Fitbauer -- nota matemática de implementación}
\date{\today}

\begin{document}
\maketitle

\section{Motivación}

Las distribuciones unidimensionales $P(\BHF)$, $P(\DEQ)$ y $P(\delta)$ son
adecuadas cuando la fuente dominante de desorden es, respectivamente, magnética,
cuadrupolar o de desplazamiento isomérico. En materiales amorfos, ferritas
desordenadas, vidrios, nanopartículas o fases con múltiples entornos locales
puede ocurrir que dos parámetros cambien a la vez. En ese caso se puede plantear
una distribución bidimensional:
\begin{equation}
  P(x,y)\ge 0,\qquad x,y\in\{\BHF,\DEQ,\delta\}.
\end{equation}
Los casos implementados en la GUI son $P(\BHF,\DEQ)$, $P(\delta,\DEQ)$ y
$P(\BHF,\delta)$. El interés principal es poder estudiar no sólo las marginales,
sino también la posible correlación entre las dos magnitudes distribuidas. El
caso $P(\delta,\DEQ)$ es especialmente relevante para dobletes paramagnéticos
anchos en vidrios, silicatos y fases amorfas.

\section{Modelo espectral}

Sea $v_k$ el canal de velocidad y sea $S(v;B,Q,\delta,\Gamma)$ el sextete de
absorción de profundidad unitaria. En una malla $x_i$, $y_j$ el modelo lineal en
los pesos es
\begin{equation}
  T_k = b_0 + b_1 v_k - \sum_{i=1}^{N_x}\sum_{j=1}^{N_y}
  P_{ij}\,S(v_k;B_{ij},Q_{ij},\delta_{ij},\Gamma).
\end{equation}
Los parámetros no distribuidos se mantienen globales. Por ejemplo, en
$P(\delta,\DEQ)$ se fija o ajusta un $B$ global; si $B=0$ el patrón representa un
conjunto de dobletes colapsados magnéticamente.
En forma matricial,
\begin{equation}
  \bm y \simeq X\bm z,
\end{equation}
donde $\bm z$ contiene fondo, pendiente y los pesos $P_{ij}$ aplanados. Las
columnas asociadas a la distribución son $-S(v_k;B_i,Q_j)$ porque la absorción
resta transmisión.

\section{Regularización 2D}

El problema inverso es fuertemente subdeterminado. Por ello se penaliza la
curvatura de la distribución en las dos direcciones:
\begin{equation}
  \Phi(\bm z)=\|W(\bm y-X\bm z)\|^2
  +\alpha_B\|D_B^2 P\|^2
  +\alpha_Q\|D_Q^2 P\|^2,
\end{equation}
con la restricción
\begin{equation}
  P_{ij}\ge 0.
\end{equation}
Aquí $W$ es el peso estadístico de los canales, $D_B^2$ aplica segundas
diferencias a lo largo de $\BHF$, y $D_Q^2$ aplica segundas diferencias a lo
largo de $\DEQ$. En la implementación inicial se usa una malla rectangular y los
operadores se escriben mediante productos de Kronecker:
\begin{align}
  L_B &= D_B^2\otimes I_Q,\\
  L_Q &= I_B\otimes D_Q^2.
\end{align}
El problema aumentado resuelto por mínimos cuadrados acotados es
\begin{equation}
  \begin{bmatrix}
    WX\\
    \sqrt{\alpha_B}\,L_B\\
    \sqrt{\alpha_Q}\,L_Q
  \end{bmatrix}\bm z
  \simeq
  \begin{bmatrix}
    W\bm y\\
    \bm 0\\
    \bm 0
  \end{bmatrix},
  \qquad P_{ij}\ge 0.
\end{equation}

\section{Marginales, momentos y correlación}

Una vez ajustada $P_{ij}$, las distribuciones marginales son
\begin{align}
  P_B(B_i) &= \int P(B_i,Q)\,\dd Q \simeq \sum_j P_{ij}\,\Delta Q_j,\\
  P_Q(Q_j) &= \int P(B,Q_j)\,\dd B \simeq \sum_i P_{ij}\,\Delta B_i.
\end{align}
En una malla no uniforme deben usarse pesos de integración adecuados. La
probabilidad bidimensional normalizada se define por
\begin{equation}
  p_{ij}=\frac{P_{ij}}{\int\!\int P(B,Q)\,\dd B\,\dd Q}.
\end{equation}
Con $p_{ij}$ se calculan los momentos diagnósticos
\begin{align}
  \langle B\rangle &= \sum_{ij} p_{ij} B_i\,\Delta B_i\Delta Q_j, &
  \langle Q\rangle &= \sum_{ij} p_{ij} Q_j\,\Delta B_i\Delta Q_j,\\
  \sigma_B^2 &= \sum_{ij} p_{ij}(B_i-\langle B\rangle)^2\Delta B_i\Delta Q_j, &
  \sigma_Q^2 &= \sum_{ij} p_{ij}(Q_j-\langle Q\rangle)^2\Delta B_i\Delta Q_j,
\end{align}
y la correlación aparente
\begin{equation}
  \rho_{BQ}=\frac{\sum_{ij}p_{ij}(B_i-\langle B\rangle)(Q_j-\langle Q\rangle)
  \Delta B_i\Delta Q_j}{\sigma_B\sigma_Q}.
\end{equation}
Una nube inclinada en el mapa $B$--$Q$ puede sugerir correlación entre campo y
cuadrupolo, pero esa interpretación sólo es válida si la solución es estable.

\section{Componentes nítidas simultáneas}

La implementación permite añadir componentes nítidas al problema 2D. El modelo
se extiende a
\begin{equation}
  T_k=b_0+b_1v_k-\sum_{ij}P_{ij}S_{ij}(v_k)-\sum_m a_m C_m(v_k)-F_k,
\end{equation}
donde $C_m$ son componentes discretas con amplitud libre $a_m\ge0$ y $F_k$ es
la absorción conocida de componentes con profundidad fijada. Las columnas nítidas
no se regularizan; compiten con la distribución sólo a través del residuo. Esto
es útil para separar una fase discreta conocida de un fondo desordenado, pero
también puede transferir peso artificialmente entre mapa y nítidos si solapan.

\section{L-surface de regularización 2D}

En 1D se usa una curva L. En 2D hay dos regularizaciones, por lo que el análogo
es una superficie en $(\alpha_B,\alpha_Q)$. Para cada pareja se calculan
\begin{align}
  R(\alpha_B,\alpha_Q)&=\|W(\bm y-X\bm z)\|,\\
  G(\alpha_B,\alpha_Q)&=\sqrt{\|D_B^2P\|^2+\|D_Q^2P\|^2},\\
  \mathrm{RMS}(\alpha_B,\alpha_Q)&=\sqrt{\langle r_k^2\rangle}.
\end{align}
La GUI muestra un mapa de RMS frente a $\log_{10}\alpha_B$ y
$\log_{10}\alpha_Q$ y una proyección $R$ frente a $G$. Se proponen dos puntos:
el mínimo de RMS y un compromiso geométrico entre residuo y rugosidad. Ninguno
es una respuesta automática; debe verificarse visualmente la estabilidad del
mapa.

\section{Riesgo de sobreajuste}

\textbf{Advertencia esencial:} una distribución 2D puede ajustar aparentemente
muy bien y no tener significado físico. Por ejemplo, una malla $30\times30$
contiene 900 pesos, a menudo más que el número efectivo de puntos independientes
del espectro. Con regularización débil, el mapa puede absorber:
\begin{itemize}
  \item ruido estadístico,
  \item errores de folding o de escala de velocidad,
  \item anchuras de línea mal elegidas,
  \item fondo instrumental,
  \item componentes discretas omitidas,
  \item relajación dinámica que debería modelarse de otro modo.
\end{itemize}
Por tanto, un residuo bajo o un mapa visualmente atractivo no demuestran por sí
solos la existencia de una distribución física $P(\BHF,\DEQ)$.

\section{Buenas prácticas}

Antes de interpretar el mapa se recomienda:
\begin{enumerate}
  \item comparar con modelos más simples: sextetes discretos, $P(\BHF)$,
  $P(\DEQ)$ y componentes nítidas;
  \item variar $\alpha_B$ y $\alpha_Q$ en varios órdenes de magnitud;
  \item repetir con distinto número de bins;
  \item verificar que las marginales y rasgos principales son estables;
  \item fijar $\delta$ y $\Gamma$ si se conocen independientemente;
  \item comprobar que el residuo no conserva estructura sistemática;
  \item contrastar con información externa: TEM, magnetometría, estructura,
  química o literatura.
\end{enumerate}
Para publicación, el modelo 2D debería justificarse mostrando que un modelo 1D o
discreto no explica los datos de forma satisfactoria y que la distribución 2D es
robusta frente a cambios razonables de regularización.

\section{Estado de implementación}

La implementación de Fitbauer incluye el backend:
\begin{center}
\texttt{mossbauer\_distribution.fit\_bhf\_quad\_distribution()}
\end{center}
y una integración en la GUI como modos \textbf{P(BHF,$\DEQ$) 2D},
\textbf{P($\delta$,$\DEQ$) 2D}, \textbf{P(BHF,$\delta$) 2D} y como distribución
1D \textbf{P($\delta$)}. El backend devuelve centros de los dos ejes, matriz de
pesos, probabilidad normalizada, curva ajustada, residuo, fondo, pendiente,
valores de $\alpha_x$, $\alpha_y$, marginales, momentos, correlación aparente y
avisos de identificabilidad. La GUI muestra un mapa de calor con marginales,
permite exportar el mapa a TSV, incluye una L-surface $\alpha_x$--$\alpha_y$,
dibuja el heatmap en Plotly y añade el mapa 2D al PDF del informe. También se
admiten componentes nítidas simultáneas. Quedan para fases posteriores mejoras
de rendimiento para mallas muy grandes y validaciones avanzadas específicas por
tipo de muestra.

\end{document}
